\chapter{Eine Solver-Architektur für prüfbare Kandidaten}

\section{Architektur beginnt an einer Verantwortungsgrenze}

\paragraph{Lesebrücke.}
Das vorangehende Kapitel hat einen Kandidatenvertrag erzeugt. Architektur ist
nötig, weil dieser Vertrag verloren gehen kann, wenn Parsen, Lösen,
Visualisieren und Modellbearbeitung als eine Operation behandelt werden.

Ein Solver-Dienst darf eine schreibgeschützte Problemerklärung einlesen und
Kandidatenevidenz zurückgeben. Er darf weder den ingenieurtechnischen Zweck
ableiten noch die freigegebenen Freiheitsgrade erweitern, Quellbeobachtungen
ersetzen, Anwendbarkeit entscheiden oder ein Ergebnis allein wegen des
Abbruchs eines numerischen Laufs in ein autoritatives Alignment schreiben.

Für die durchgehende Abweichung lautet der konzeptionelle Ablauf
\[
\begin{array}{c}
\boxed{\mathcal P}\longrightarrow\boxed{\text{prüfen und zusammenstellen}}\\[0.4em]
\downarrow\\[-0.2em]
\boxed{\text{lösen und verifizieren}}\longrightarrow\boxed{\mathcal C}\\[0.4em]
\downarrow\\[-0.2em]
\boxed{\text{evaluieren und entscheiden}}
\end{array}
\]
Nur die mittleren Operationen gehören zur Solver-Verantwortung. Der letzte
Pfeil ist eine Grenze zum ingenieurtechnischen Ablauf.

\section{Das Problempaket}

Ein übertragbares Problempaket benötigt mehr als Geometrie:
\[
\mathcal P=
(id,I,P,M,x_0,\mathcal X_{\mathrm{free}},q,\mathcal F,
J,W,\Sigma,\pi),
\]
wobei \(W\) die Skalierung und \(\pi\) die Provenienz dokumentiert. Die
Prüfung stellt sicher, dass jede freie Variable erklärt ist, feste Eingaben
gesperrt bleiben, Einheiten zusammenpassen, Zielterme ihre Quelle nennen und
alle referenzierten Modelle und Kontexte vorliegen. Ein Scheitern dieser
Prüfung ist keine Solver-Nichtkonvergenz, sondern eine unvollständige oder
nicht anwendbare Problemerklärung.

\section{Die Zusammenstellung macht mathematische Struktur sichtbar}

\paragraph{Lesebrücke.}
Ist der ingenieurtechnische Vertrag vollständig, darf die Architektur ihn in
numerische Objekte übersetzen, ohne seine Bedeutung zu verändern.

Die Zusammenstellung erzeugt:
\begin{itemize}
\item Parametervektor und Grenzen aus den freigegebenen Freiheitsgraden;
\item Gleichungs- und Ungleichungsfunktionen aus erklärten Nebenbedingungen;
\item getrennt typisierte Ziel- und Diagnoseresiduen;
\item Rekonstruktions- und Beobachtungsevaluatoren;
\item Skalierung, Ableitungsanbieter und Muster dünner Besetzung;
\item Abbruchtoleranzen und Einstellungen für unabhängige Verifikation.
\end{itemize}

Residuen werden nicht austauschbar, nur weil sie in einem Vektor gespeichert
sind. Typ und Provenienz müssen die Zusammenstellung überdauern, damit eine
Prüfung Endpunktbedingung, Beobachtungsabweichung, Regularisierung und
numerische Diagnose unterscheiden kann.

\section{Warum SQP, Ableitungen und dünne Besetzung hinter Schnittstellen gehören}

\paragraph{Lesebrücke.}
SQP ist für ein glattes nichtlineares Problem mit sichtbaren Gleichungs- und
Ungleichungsnebenbedingungen nützlich. Es ist ein austauschbares numerisches
Verfahren, nicht Eigentümer des Problems.

Bei \(x_k\) bildet eine SQP-Implementierung ein lokales quadratisches Modell
und linearisierte Nebenbedingungen, bestimmt einen Schritt \(p_k\) und wendet
eine Globalisierungsregel wie Liniensuche oder Vertrauensbereich an. Die
Architektur muss
\[
x_k,\quad J(x_k),\quad r(x_k),\quad
\nabla J(x_k),\quad A(x_k),\quad p_k
\]
sowie Nebenbedingungsabstände, Ableitungsprovenienz, Skalierung und
Solver-Diagnosen offenlegen.

Analytische, automatische oder finite Differenzen dürfen Ableitungen liefern,
wenn Verfahren und Validierung dokumentiert sind. Dünne Besetzung ist eine
rechnerische Eigenschaft des erklärten Abhängigkeitsgraphen. Ihre Nutzung kann
Aufwand verringern; sie rechtfertigt kein unerklärtes Ziel und keine
unerklärte Nebenbedingung.

\section{Abbruch, Konvergenz und Verifikation}

Ein Solver kann enden, weil eine Toleranz erreicht ist, der Fortschritt
stockt, eine Iterationsgrenze erreicht oder eine Auswertung fehlgeschlagen ist.
Diese Ergebnisse müssen getrennt bleiben. Ein nominelles
Konvergenzkennzeichen ist nur mit der zugehörigen Evidenz belastbar:
\begin{itemize}
\item Ziel- und Schrittverläufe;
\item Gleichungsresiduen und Abstände zu Ungleichungsgrenzen;
\item ein zum Verfahren passendes Stationaritäts- oder Optimalitätsmaß;
\item gegebenenfalls Ableitungs- und Skalierungsprüfungen;
\item Rekonstruktion auf dem erklärten Verifikationsgitter;
\item Warnungen, aktive Annahmen und numerische Fehlerschätzungen.
\end{itemize}

Diese Evidenz belegt, was innerhalb von \(M\) und \(\mathcal P\) geschah. Sie
belegt weder die Richtigkeit des Beobachtungsmodells noch die Baubarkeit des
Kandidaten oder dass sein Zweck die Übernahme rechtfertigt.

\section{Kandidatenpaket und Prüfgrenze}

Der Dienst gibt, ohne anzuwenden, zurück:
\[
\mathcal C=
(\mathcal P.id,M.id,x^\star,\widehat\gamma^\star,
r^\star,J^\star,\text{Abstände},\text{Abbruch},
\text{Verifikation},\Sigma,\pi).
\]
Eine Prüfoberfläche darf Kandidaten vergleichen, Residuen und
Empfindlichkeiten untersuchen und eine Darstellung als Vorschau zeigen. Eine
Vorschau ist keine Physical Realization; ein Vergleich ist keine Evaluation
unter jedem relevanten Zweck. Eine »Anwenden«-Operation würde in einem
künftigen Ablauf einen getrennten autorisierten Befehl mit erhaltener
Provenienz erfordern.

\section{Aktuelle AXTRAN-Evidenz und allgemeine Architektur}

Die vorstehende Architektur ist ein wissenschaftliches
Verantwortungsmodell und eine Implementierungsrichtung. Die aktuelle
Repository-Evidenz für AXTRAN ist enger: AXTRAN ist experimentell; ein
einzelner SQP-Schritt und eine endlich bleibende wiederholte Schleife belegen
einen numerischen Pfad. Sie belegen keine allgemeine Optimiererarchitektur,
Produktionsrobustheit, Live-, Netz-, Fahrzeug- oder
Überhöhungsoptimierung und keine automatische Modellintegration.

Beispiele reichhaltiger Sitzungen, Visualisierung, alternativer Solver oder
Netzzusammenstellung sind daher künftige oder problemspezifische
Möglichkeiten. Sie dürfen nicht als gegenwärtige AIM- oder AXTRAN-Fähigkeit
gelesen werden.

\section{Was die Ingenieurin oder der Ingenieur erhält}

Für die CAD-/Vermessungsabweichung gibt die Architektur einen oder mehrere
nachvollziehbare Krümmungskandidaten und die für ihre Befragung erforderliche
Evidenz zurück: Welche Parameter änderten sich, was blieb fest, welche
Residuen verbesserten sich, welche Nebenbedingungen sind nahezu aktiv, warum
endete der Lauf und welche Unsicherheiten bleiben? Danach wird die
Anwendbarkeit im vorgesehenen Projekt- und Realisierungskontext geprüft.
Evaluation vergleicht Folgen unter erklärten Zwecken. Eine autorisierte
Engineering Decision darf annehmen, ablehnen oder eine neue
Problemformulierung verlangen.

\begin{summary}
Eine Solver-Architektur ist vertrauenswürdig, wenn sie das
ingenieurtechnische Problem beim Eintritt bewahrt und beim Austritt einen
prüfbaren Kandidaten liefert.

\[
\boxed{\text{berechnet}\neq\text{anwendbar}\neq
\text{angenommen}\neq\text{autoritativ}}
\]

Der Solver verantwortet numerische Arbeit innerhalb des erklärten Vertrags.
Er besitzt weder Alignment-Identität noch Anwendbarkeit, Annahme oder
Autorität.
\end{summary}
